\chapter{Experiments and Evaluation}

This chapter presents the evaluation setup, metrics, results, and interpretation. The reported values come from the JSON quality report generated by the project quality-suite script.

\section{Evaluation Setup}

The evaluation uses a GTU-specific question set stored as JSON. Each question includes expected source indicators and expected answer keywords. The cases cover regulations, guides, web pages, department information, and PDF documents. The quality suite submits each question through the same manual-question path used by the application, then inspects the stored answer and retrieval traces.

The recorded run used a local SQLite quality-suite database. This local mode is useful for repeatable development tests. Production deployment can use PostgreSQL with pgvector for vector operations and database extensions.

\section{Evaluation Questions}

The 12 questions were grouped into four categories:

\begin{itemize}
    \item \textbf{Regulations:} undergraduate semester duration, graduate education regulation, thesis writing guide, diploma directive.
    \item \textbf{Web pages:} ECTS package, academic calendar, regulation page, Computer Engineering department description.
    \item \textbf{Department information:} Computer Engineering mission and student profile.
    \item \textbf{PDF documents:} research policy, Erasmus directive, 2025--2026 student guide.
\end{itemize}

\section{Metrics}

The quality suite calculates five primary metrics. Source hit rate measures whether any retrieved trace matches the expected source hints. Top-1 source hit rate is stricter and only considers the first retrieved trace. Answer keyword hit rate checks whether the generated response includes expected key terms. Fallback rate measures whether the system used deterministic fallback output instead of the provider model. Latency measures end-to-end answer generation time.

\section{Results}

Table \ref{tab:evaluation-summary} shows the summary from the latest recorded evaluation.

\begin{table}[!htbp]
    \centering
    \caption{\label{tab:evaluation-summary}Evaluation summary for the 12-question GTU set.}
    \begin{tabular}{lr}
        \toprule
        Metric & Value \\
        \midrule
        Question count & 12 \\
        Top-1 source hit rate & 83.33\% \\
        Source hit rate & 100.00\% \\
        Answer keyword hit rate & 91.67\% \\
        Fallback rate & 0.00\% \\
        Average latency & 7.75 s \\
        Median latency & 6.58 s \\
        P95 latency & 10.60 s \\
        Maximum latency & 15.79 s \\
        Average confidence & 0.5931 \\
        \bottomrule
    \end{tabular}
\end{table}

\section{Per-Category Observations}

Regulation and directive questions performed well when the expected source was a distinctive PDF. Examples include the undergraduate regulation, thesis writing guide, diploma directive, and Erasmus directive. These questions often contain regulation or guide terms that the retrieval system can match in filenames, titles, or PDF content.

Some web navigation questions were harder. The academic-calendar and lisans regulation-listing questions had weaker top-1 behavior because other official sources contained related terms but were not the most precise answer source: a student handbook PDF and the undergraduate regulation PDF respectively ranked first. This shows that official but indirect sources can outrank direct pages when the corpus contains broad student guides or regulation documents.

The Computer Engineering department questions retrieved the correct department page in every case. For one of them, however, the passage selected for the model did not include the exact wording about training creative and productive engineers, so the generated answer reported that the information was not available. This is a passage-selection precision gap rather than a retrieval failure.

\section{Latency Analysis}

The median latency was 6.58 seconds and P95 latency was 10.60 seconds. This is acceptable for a controlled demo but high for a polished live production experience. The latency includes retrieval, provider calls, database writes, and answer post-processing. If TTS generation is included in the viewer experience, perceived latency can be higher unless audio generation is decoupled from screen answer display.

The largest observed latency was 15.79 seconds for the academic-calendar web question. This suggests that navigation-style questions may need better direct-page routing or cached canonical answers.

\section{Failure and Weakness Analysis}

The system did not use fallback output in the recorded run. One answer did not contain the exact expected keywords: for the Computer Engineering purpose question, the relevant wording was not present in the passage given to the model, so the model reported missing information rather than quoting the official phrasing. Top-1 source precision is the other main weakness. When the first source is wrong but a correct source appears lower in the retrieved traces, the generated answer may still be acceptable, but debugging and trust are weaker. For a university-facing assistant, improving first-ranked evidence is important.

The top-1 misses were:

\begin{itemize}
    \item academic calendar question: the correct calendar page appeared in the traces, but a student handbook PDF ranked first;
    \item lisans regulation-listing question: the correct listing page appeared in the traces, but the undergraduate regulation PDF (YN-0001) ranked first.
\end{itemize}

These misses indicate that the system would benefit from stronger general signals for canonical pages, such as down-weighting broad PDFs for navigation-style questions and better PDF title normalization, rather than per-question hard-coded routing that would not generalize. The keyword miss also points to passage selection: the most relevant sentence is not always the one surfaced to the model.

\section{Automated Tests}

The backend includes automated tests for ingestion and manual question flow, manual queue live-state behavior, live-stage playback timing, ambient segment behavior, evaluation metric calculations, provider option resolution, LLM JSON parsing, and TTS runtime settings. These tests protect core behavior in the prototype and are especially useful because many production paths can be exercised without a live provider key.

\section{Discussion}

The results show that the project reached the goal of a working controlled demo. The answer coverage is strong, the source hit rate is usable, and the system can operate with real application services rather than isolated scripts. At the same time, the evaluation makes the next engineering priorities clear: improve top-1 retrieval precision, reduce answer latency, expand the official GTU corpus, and add stronger monitoring around weak-context answers.
